\documentclass[a4paper, serif, 10pt]{moderncv}

%% moderncv
% color and style
\moderncvcolor{black}
\moderncvstyle{banking}

% renew moderncv command to adapt for CJK colon
\renewcommand*{\cvitem}[3][.25em]{%
  \ifstrempty{#2}{}{\hintstyle{#2}：}{#3}%
  \par\addvspace{#1}}

\renewcommand*{\cvitemwithcomment}[4][.25em]{%
  \savebox{\cvitemwithcommentmainbox}{\ifstrempty{#2}{}{\hintstyle{#2}：}#3}%
  \setlength{\cvitemwithcommentmainlength}{\widthof{\usebox{\cvitemwithcommentmainbox}}}%
  \setlength{\cvitemwithcommentcommentlength}{\maincolumnwidth-\separatorcolumnwidth-\cvitemwithcommentmainlength}%
  \begin{minipage}[t]{\cvitemwithcommentmainlength}\usebox{\cvitemwithcommentmainbox}\end{minipage}%
  \hfill% fill of \separatorcolumnwidth
  \begin{minipage}[t]{\cvitemwithcommentcommentlength}\raggedleft\small\itshape#4\end{minipage}%
  \par\addvspace{#1}}

%% page layout/margins
\usepackage[top=2.5cm, bottom=2.5cm, left=1.5cm, right=1.5cm]{geometry}

\nopagenumbers{}

%% Babel config for English language
\usepackage[english]{babel}

%% fontspec
\usepackage{fontspec}

\IfFontExistsTF{Linux Libertine}{
  \setmainfont[Ligatures={TeX, Common}, Numbers=Lining, ItalicFont=Linux Libertine]{Linux Libertine}
}{}
\IfFontExistsTF{Linux Libertine O}{
  \setmainfont[Ligatures={TeX, Common}, Numbers=Lining, ItalicFont=Linux Libertine O]{Linux Libertine O}
}{}

%% CTeX
% CJK support, used to show CJK characters in the resume
%
% - heading=false: disable ctex heading
% - punct=kaiming: use kaiming punctuations styles for CJK
% - scheme=plain: use plain scheme, do not override \normalsize font size
% - space=auto: space settings for CJK characters
%
% CJK fonts are resolved per locale language, e.g. Noto Serif/Sans CJK JP
% for Japanese resumes, with fallback to the next region if not installed
%
% ref:
% - http://ctan.mirrorcatalogs.com/language/chinese/ctex/ctex.pdf
\usepackage[UTF8, heading=false, punct=kaiming, scheme=plain, space=auto]{ctex}

\IfFontExistsTF{Noto Serif CJK HK}{
  \setCJKmainfont{Noto Serif CJK HK}
}{
\IfFontExistsTF{Noto Serif CJK TC}{
  \setCJKmainfont{Noto Serif CJK TC}
}{}}
\IfFontExistsTF{Noto Sans CJK HK}{
  \setCJKsansfont{Noto Sans CJK HK}
}{
\IfFontExistsTF{Noto Sans CJK TC}{
  \setCJKsansfont{Noto Sans CJK TC}
}{}}

%% line spacing
\usepackage{setspace}
\setstretch{1.125}

%% URL styling - use normal text instead of monospace
\urlstyle{same}

%% Patch \httplink and \httpslink to handle full URLs
% The original moderncv macros always prepend http:// or https:// to URLs.
% This causes issues when the URL already has a protocol (e.g., https://example.com)
% resulting in malformed URLs like https://https://example.com.
% This patch checks if the URL already contains "://" and uses it directly if so.
% Note: We use \str_set:Nx (x-expansion) to expand macros like \@homepage first.
\ExplSyntaxOn
\str_new:N \g_yamlresume_url_str

\RenewDocumentCommand{\httpslink}{O{}m}{
  \str_set:Nx \g_yamlresume_url_str {#2}
  \str_if_in:NnTF \g_yamlresume_url_str {://}
    { \url{#2} }
    { \url{https://#2} }
}

\RenewDocumentCommand{\httplink}{O{}m}{
  \str_set:Nx \g_yamlresume_url_str {#2}
  \str_if_in:NnTF \g_yamlresume_url_str {://}
    { \url{#2} }
    { \url{http://#2} }
}
\ExplSyntaxOff

\name{陳偉明}{}
\title{資深業務分析師 | 數據驅動決策與流程優化}
\phone[mobile]{+852 9123 4567}
\email{chan.waiming@example.com}
\homepage{https://www.linkedin.com/in/chanwaiming}

\address{中國香港，香港島，香港，中環皇后大道中99號}{}{}

\extrainfo{{\small \faLinkedin}\ \href{https://www.linkedin.com/in/chanwaiming}{@chanwaiming}}

\begin{document}

\maketitle

\section{簡介}

\cvline{}{\begin{itemize}
\item 擁有超過6年業務分析經驗，專注於金融與電訊行業的數字化轉型
\item 擅長需求收集、業務流程建模、數據分析及利益相關者管理
\item 精通敏捷開發方法，熟練運用SQL、Tableau、Power BI等工具
\item 具備優秀的溝通與協調能力，能有效橋接業務與技術團隊
\item 持有CBAP、PMP及IIBA-AAC等專業認證
\end{itemize}}
\section{教育背景}

\cventry{2014年9月–2018年6月}
        {學士，工商管理學士（資訊系統），成績：3.6/4.0}
        {香港大學}
        {\url{https://www.hku.hk}}
        {}
        {\begin{itemize}
\item 主修資訊系統，副修統計學，成績優異
\item 曾獲院長嘉許名單，並於畢業論文研究電子商務用戶行為
\item 積極參與商業案例競賽，獲得校內二等獎
\end{itemize}
\textbf{課程}：商業分析、資料庫管理、系統分析與設計、項目管理、統計學、財務會計、市場營銷、運營管理}
\section{工作經歷}

\cventry{2021年3月至今}
        {高級業務分析師}
        {匯豐銀行}
        {\url{https://www.hsbc.com.hk}}
        {}
        {\begin{itemize}
\item 領導零售銀行數字化轉型項目，負責需求收集與業務流程分析
\item 與IT團隊協作，將業務需求轉化為技術規格，確保交付符合預期
\item 主導CRM系統升級，提升客戶數據整合效率達30\%
\item 建立實時業務儀表板，監控關鍵績效指標，輔助管理層決策
\item 指導初級分析師，推動團隊採用敏捷實踐
\end{itemize}
\textbf{關鍵字}：需求收集、流程改善、數據分析、敏捷開發、SQL、Tableau、利益相關者管理}

\cventry{2018年7月–2021年2月}
        {業務分析師}
        {香港電訊}
        {\url{https://www.hkt.com}}
        {}
        {\begin{itemize}
\item 負責電訊業務支撐系統的需求分析與流程優化
\item 使用BPMN和UML繪製業務流程圖，識別瓶頸並提出改善方案
\item 撰寫用戶故事與驗收標準，協調UAT測試，確保系統上線質量
\item 參與跨部門項目，成功將訂單處理時間縮短20\%
\end{itemize}
\textbf{關鍵字}：業務流程建模、BPMN、UML、用戶故事、UAT、流程優化}

\cventry{2017年6月–2017年8月}
        {業務分析實習生}
        {香港交易所}
        {\url{https://www.hkex.com.hk}}
        {}
        {\begin{itemize}
\item 協助分析市場數據，製作每週交易報告
\item 參與需求研討會，記錄並整理業務需求
\item 支援團隊完成系統測試案例編寫
\end{itemize}
\textbf{關鍵字}：數據分析、報告製作、需求整理}
\section{語言}

\cvline{粵語}{母語或雙語水平 \hfill \textbf{關鍵字}：母語}
\cvline{英語}{完全專業水平 \hfill \textbf{關鍵字}：IELTS 7.5、商業書寫}
\cvline{普通話}{有限工作水平 \hfill \textbf{關鍵字}：HSK 5、基本會話}
\section{專業技能}

\cvline{業務分析}{專家 \hfill \textbf{關鍵字}：需求收集、流程改善、差距分析、利益相關者管理、業務流程建模、用戶故事、驗收測試}
\cvline{數據分析}{高級 \hfill \textbf{關鍵字}：SQL、Python、Tableau、Power BI、Excel、統計分析、A/B測試}
\cvline{項目管理}{高級 \hfill \textbf{關鍵字}：敏捷開發、Scrum、瀑布模型、JIRA、Confluence、風險管理}
\cvline{工具}{中級 \hfill \textbf{關鍵字}：Visio、Lucidchart、UML、BPMN、MS Project、Google Analytics}
\section{榮譽}

\cventry{2023年12月}
        {匯豐銀行}
        {傑出員工獎}
        {}
        {}
        {表彰在數字化轉型項目中的卓越貢獻，成功提升客戶數據整合效率。}

\cventry{2017年6月}
        {香港大學}
        {院長嘉許名單}
        {}
        {}
        {表揚學術成績優異，位列學院前10\%。}
\section{證書}

\cventry{2022年11月}
        {IIBA}
        {認證業務分析師 (CBAP)}
        {\url{https://www.iiba.org/certification/cbap/}}
        {}
        {}

\cventry{2021年5月}
        {IIBA}
        {敏捷分析認證 (IIBA-AAC)}
        {\url{https://www.iiba.org/certification/agile-analysis/}}
        {}
        {}

\cventry{2020年8月}
        {PMI}
        {Project Management Professional (PMP)}
        {\url{https://www.pmi.org/certifications/project-management-pmp}}
        {}
        {}
\section{出版刊物}

\cventry{2023年3月}
        {數據驅動的業務分析：香港金融業的實踐}
        {香港資訊科技雜誌}
        {\url{https://www.hkitjournal.com/data-driven-business-analysis}}
        {}
        {探討如何利用數據分析提升業務決策效率，並分享香港金融業的實際案例。}
\section{項目}

\cventry{2023年1月–2023年12月}
        {為匯豐銀行零售銀行部門升級CRM系統，整合客戶數據並提升銷售效率。}
        {客戶關係管理系統升級}
        {}
        {}
        {\begin{itemize}
\item 負責需求收集與業務流程分析，定義系統功能與用戶故事
\item 協調IT團隊與業務部門，確保系統按時交付
\item 設計UAT測試案例，培訓用戶，系統上線後客戶滿意度提升15\%
\end{itemize}
\textbf{關鍵字}：CRM、Salesforce、需求收集、UAT}

\cventry{2022年3月–2022年9月}
        {為香港電訊建立實時業務績效儀表板，監控關鍵指標。}
        {實時業務儀表板開發}
        {}
        {}
        {\begin{itemize}
\item 使用Tableau和SQL開發儀表板，視覺化呈現銷售與運營數據
\item 與管理層合作定義KPI，確保儀表板滿足決策需求
\item 儀表板上線後，報告生成時間減少50\%
\end{itemize}
\textbf{關鍵字}：Tableau、SQL、數據可視化、KPI}
\section{興趣愛好}

\cvline{遠足}{龍脊、鳳凰山、麥理浩徑}
\cvline{攝影}{風景攝影、街頭攝影}
\cvline{閱讀}{商業分析、數據科學、歷史}
\section{志願服務}

\cventry{2021年9月至今}
        {職涯導師}
        {香港青年協會}
        {\url{https://www.hkfyg.org.hk}}
        {}
        {\begin{itemize}
\item 為大學生提供職涯規劃指導，分享業務分析行業經驗
\item 協助舉辦模擬面試及履歷工作坊，累計服務超過50名學生
\end{itemize}}
    
\end{document}